
Existing efficient MoE kernels also frame MoE computation as a Grouped GEMM, but their ingredients are different from \Algname. Here we provide an overview (but not a complete list) of key differences:

\begin{itemize}
    \item \textbf{ScatterMoE} \citep{tan2024scattered}\footnote{\scriptsize \url{https://github.com/shawntan/scattermoe/blob/47b5e1502e5a10e82c8e5945d761b877849871e7/scattermoe/mlp.py\#L51}} implements gather fusion for varlen-$\mathbf{M}$ Grouped GEMM but not for varlen-$\mathbf{K}$ Grouped GEMM. ScatterMoE also does not overlap MMA computation with memory IO. Moreover, ScatterMoE is also built on older versions of Triton where TMA is not supported. ScatterMoE computes $dS$ as $dS = \langle dO, Y \rangle$ which requires caching $Y$. This results in large IO cost and activation memory requirement. Both ScatterMoE's forward and backward pass have limited fusion and hence ScatterMoE is much slower than~\Algname, especially for the backward pass.

    \item \textbf{MoMoE} \citep{costin2025momoe}\footnote{\scriptsize \url{https://github.com/tilde-research/MoMoE-impl/blob/d6e2d683185bfe4030265c3ca062564356faa61e/momoe/momoe.py\#L914}} also implements the gather fusion for varlen-$\mathbf{M}$ but not varlen-$\mathbf{K}$ Grouped GEMM similarly to ScatterMoE. Although fused with up-proj activation gradient, the $dS$ computation still utilizes $dS = \langle dO, Y \rangle$. Similar to ScatterMoE, MoMoE does not use TMA for IO. The scatter operation in MoMoE is (much) slower than~\Algname, as shown in Figure~\ref{fig:scatter-and-add-speed}.
    
    \item \textbf{MegaBlocks} \citep{gale2023megablocks} has multiple MoE implementations and we focus on $\mathrm{ParallelDroplessMLP}$\footnote{\scriptsize \url{https://github.com/databricks/megablocks/blob/78eea65fda01e638af36ae38853bc51efb04a4b4/megablocks/layers/dmoe.py\#L18}} which is built on top of block-sparse matrix multiplication\footnote{\scriptsize\url{https://github.com/databricks/megablocks/blob/78eea65fda01e638af36ae38853bc51efb04a4b4/megablocks/layers/mlp.py\#L308}}. $\mathrm{ParallelDroplessMLP}$ first gathers and pads the tokens and then launches block-sparse GEMM for up and down-proj. Then, it launches a scatter kernel before reducing across the expert results. These sparse matrix multiplications usually take a longer time than the highly-optimized Grouped GEMM, as shown in Figure~\ref{fig:moe_breakdown}, and the gather and scatter kernel have a total IO cost of $8TKd$ bytes which can be a bottleneck for fine-grained MoEs. We consider MegaBlocks's $\mathrm{ParallelDroplessMLP}$ as a block-sparse GEMM baseline in our benchmark and find that MoE implemented via Grouped GEMMs often have a higher training throughput than MoEs implemented via block-sparse GEMM. 

    \item \textbf{Megatron-LM} \citep{shoeybi2019megatron} also has multiple MoE implementations and we focus on $\mathrm{GroupedMLP}$\footnote{\scriptsize\url{https://github.com/NVIDIA/Megatron-LM/blob/610a75ef3a4a80c2ce2da436c19244e5362978d4/megatron/core/transformer/moe/experts.py\#L100}},  which uses Grouped GEMM\footnote{\scriptsize \url{https://github.com/fanshiqing/grouped_gemm/blob/main/csrc/grouped_gemm.cu}} from the CUTLASS library \citep{CUTLASS_gemm} with JIT epilogue fusion as the GEMM backend. Similar to DeepGEMM, $\mathrm{GroupedMLP}$ does not fuse gather with the prologue (it assumes contiguously-packed inputs). A recent memory-efficient patch\footnote{\scriptsize \url{https://github.com/NVIDIA/Megatron-LM/commit/2659630721ac87237c8cb772b1c2f1b34176f443}} fuses $S$ weighting with SwiGLU computation during forward, and during backward which allows the PyTorch autograd engine \citep{paszke2019pytorchimperativestylehighperformance} to follow a similar computational path as \Algname.% However, we note that such reliance would still be both activation memory- and computationally-inefficient if we include a bias calculation. We could fuse $d\mathrm{Bias}$ with \Algname's $dH$ kernel epilogue and do a similar partial row-sum as $dS$ ($dS$ is performing partial column sum), and such fusion is not achievable if we have a separate Grouped GEMM call. We also find $\mathrm{GroupedMLP}$'s expert aggregation is severely under-optimized (a direct $\mathrm{torch.scatter\_add}$ call) as shown in Figure~\ref{fig:moe_breakdown}. \mayank{remove all logic about bias}
    % \TD{One extra thing to mention wrt Megatron's doign S after the swiglu in the fwd is that if there's bias, you would also need to scale the bias, and it creates a very big tensor. So our design (doing S after down proj) and manually doing the bwd is still better / cleaner / more flexible. As an aside, I personally feel using autograd really limits ones ability to fuse kernels, the current way autograd is written does not take efficiency into account.}

    Megatron-LM also implements $\mathrm{TEGroupedMLP}$\footnote{\scriptsize\url{https://github.com/NVIDIA/Megatron-LM/blob/610a75ef3a4a80c2ce2da436c19244e5362978d4/megatron/core/transformer/moe/experts.py\#L746}} which launches 4 CUDA streams to execute a list of GEMM (without contiguously-packed inputs, and without a persistent tile scheduler). In this case, each expert independently launches a new GEMM kernel leading to ``bubbles'' on the CUDA streams. This leads to underutilization of the GPU resources. We empirically find that $\mathrm{TEGroupedMLP}$ runs slower than $\mathrm{GroupedMLP}$ and so we use $\mathrm{GroupedMLP}$ across all benchmarks.

    \item \textbf{DeepGEMM} \citep{deepgemm} designs a Grouped GEMM kernel for contiguously-packed inputs. They also don't implement any other fusion for both SM90 (Hopper) and SM100 (Blackwell) BF16 Grouped GEMM. DeepGEMM specializes more on distributed training with expert parallelism \citep{deepep2025}, and it is common to launch a separate all2all kernel \citep{lepikhingshard} which is then followed by a contiguous Grouped GEMM. DeepGEMM SM90 BF16 kernel also assumes that each expert receives a multiple of $\tileM$ tokens as it does not implement the TMA tensor descriptor online update during the Grouped GEMM computation. DeepGEMM's BF16 GEMM on SM90 also does not employ Ping-Pong scheduling. DeepGEMM's varlen-$\mathbf{M}$ grouped GEMM usually uses tile shape (128, 256, 64) for both Hopper and Blackwell GPUs, while \Algname~often picks (256, 256, 64) with 2-CTA MMA for Blackwell GPUs. DeepGEMM does not use CLC persistent tile scheduler for grouped GEMM on Blackwell GPUs while \Algname~often adopts CLC persistent tile scheduler on Blackwell. 
\end{itemize}

Additionally, ScatterMoE and MoMoE are both implemented in Triton \citep{triton} for the ease of development at the expense of losing full programmability of the asynchronous compute and memory IO of Hopper and Blackwell GPUs \citep{nvidia2022_h100,nvidia2022_gb200}. For example, they cannot implement fine-grained control of asynchronous load and store during the GEMM's epilogue. They also cannot overlap MMA with heavy epilogue operations using Ping-Pong scheduling. It becomes increasingly important to overlap IO operations in epilogue when the GEMM computations are small in size (as in the case of fine-grained MoEs) to achieve high GPU utilization.
